% Verified chain from Vandiver to the cup hypothesis.
% Exact FK locators refer to the author preprint dated 15 August 2012.
\subsection{Vandiver alone supplies the cup hypothesis}
\label{sec:vandiver-kernel}

The extra cup hypothesis in Theorem~\ref{thm:pole-general} follows from
Vandiver by combining the integral results of Fukaya--Kato and
Wake--Wang-Erickson. The resulting statement is stronger than the numerical
range obtained from Sharifi's earlier cup computations.

\begin{theorem}[K17-VANDIVER-KERNEL]\label{thm:vandiver-kernel}
Let $p$ be an odd prime. If
\[
 p\nmid h\bigl(\mathbb Q(\mu_p)^+\bigr),
\]
then the odd-character cup products span $H^2(P,\mathbb F_p)$ and
\[
 \Omega_p^*=\Omega_p.
\]
The hypothesis is equivalently
\[
 K_{4k}(\mathbb Z)\{p\}=0
 \quad\text{for every }k\ge1,
\]
or equivalently the same vanishing for $1\le k\le(p-3)/2$.
These are equivalent formulations of the sufficient hypothesis;
no converse from the kernel-field equality is asserted.
\end{theorem}

\begin{proof}
For $p\ge5$, put $F_r=\mathbb Q(\mu_{p^r})$,
$F_\infty=\bigcup_rF_r$, and
\[
 X=\varprojlim_r\operatorname{Cl}(F_r)\{p\},\qquad
 \mathfrak X=\operatorname{Gal}(M/F_\infty),
\]
where $M$ is the maximal abelian pro-$p$ extension of $F_\infty$
unramified outside $p$. Thus $X$ is a quotient of $\mathfrak X$.
Let $\omega$ be the Teichm\"uller character and
$\Lambda=\mathbb Z_p[[\operatorname{Gal}(F_\infty/F_1)]]$.

First, Vandiver at $F_1$ implies $X^+=0$. Here is a proof of the
persistence statement needed for this implication. Suppose $L/K$ is a
cyclic extension of degree $p$ ramified at precisely one finite prime,
which is totally ramified, and suppose $\operatorname{Cl}(K)\{p\}=0$.
Let $E/L$ be the Hilbert $p$-class field. It is Galois over $K$, and
$G=\operatorname{Gal}(E/K)$ is a finite $p$-group. The inertia group $I$
at the ramified prime has order $p$, since $E/L$ is unramified. Let $N$
be its normal closure in $G$. All inertia groups over that prime are
conjugate, so $E^N/K$ is everywhere unramified. A nontrivial finite
$p$-group has a quotient of order $p$; class field theory and the assumed
class-group vanishing therefore give $G/N=1$. Thus $G$ is normally
generated by a single element. Its Frattini quotient has dimension at
most one over $\mathbb F_p$, so $G$ is cyclic. Consequently $N=I=G$
and $E=L$. This proves $\operatorname{Cl}(L)\{p\}=0$.
There is no extra infinite-prime ramification for odd $p$.

Apply this argument successively to $F_{r+1}^+/F_r^+$: each extension
has degree $p$ and its only ramified finite prime is the unique prime
above $p$. All the real-layer $p$-class groups are zero. Extension and
norm for the quadratic extension $F_r/F_r^+$ identify
$\operatorname{Cl}(F_r^+)\{p\}$ with
$\operatorname{Cl}(F_r)\{p\}^+$, because their compositions are
multiplication by $2$ and $1+c$, respectively. Taking inverse limits
gives $X^+=0$.

Fix an even character $\theta\ne1,\omega^2$ and set
$\chi=\omega^{-1}\theta$. The hypotheses in
Wake--Wang-Erickson, \S2.1, are then satisfied with tame level $N=1$.
The reflected character
\[
 \theta'=\omega^2\theta^{-1}=\omega\chi^{-1}
\]
is also even and different from $1,\omega^2$; its associated odd
character is $\chi'=\chi^{-1}$. Since $X_{\theta'}=0$, their
Proposition~2.2.6 says that
\[
 \mathfrak X_{(\chi')^{-1}}^{\#}(1)
   =\mathfrak X_\chi^{\#}(1)
\]
is free of rank one over the relevant Iwasawa algebra. The operation
$\#$ inverts the Galois action, and $(1)$ applies the cyclotomic twist.
They transport the coefficient ring by automorphisms, so undoing them
preserves freeness and cyclicity. Hence $\mathfrak X_\chi$ is cyclic,
and its quotient $X_\chi$ is cyclic. In addition $X_\theta=0$.
Corollaries~5.1.1--5.1.2 now show that the cuspidal Eisenstein Hecke
algebra is Gorenstein and that Sharifi's map
\[
 \varpi:H^-_\theta/I_\theta H^-_\theta
       \longrightarrow X_\chi(1)
\]
is an integral isomorphism
\cite[\S2.1, Proposition~2.2.6, Corollaries~5.1.1--5.1.2;
preprint pp.~2--4 and~11]{WakeWangErickson2017}.

For completeness, the relevant Fukaya--Kato input has no hidden
simple-root hypothesis. Gorensteinness implies their condition
$C(\mathfrak h)$, namely rational cyclicity of the Eisenstein quotient.
Put $Q=(H^-_\theta/I_\theta H^-_\theta)/(p\text{-primary torsion})$.
Over $\mathbb Q_p$, the module $Q$ is free of rank one over
$\Lambda/(\xi)[1/p]$, so its endomorphism $1-\Upsilon\varpi$
is multiplication by a scalar. Their derivative identity and
Lemma~10.2.2 make that scalar nilpotent.
The module $Q$ is $p$-torsion free, so this endomorphism is nilpotent
on $Q$ itself. Its finite geometric series makes
$\Upsilon\varpi$ an integral automorphism. After composing $\Upsilon$ with this inverse, one obtains a splitting
of the target Iwasawa module as $Q\oplus K$. The characteristic-ideal
and Fitting-ideal argument in their \S7.1.3 makes $K$ finite.
The target has no $p$-torsion, so $K=0$. This gives their integral
Theorem~7.2.8. Their \S7.1.10 applies this
surjectivity at each finite cyclotomic level
\cite[\S\S7.1.3, 7.1.8--7.1.10, Theorem~7.2.8,
\S\S7.2.10--7.2.12 and~10.2.2--10.2.3;
preprint pp.~84--86, 89 and~107]{FukayaKato2012}.

In particular, at $r=1$ the $\theta$-component of
\[
 H^2(\operatorname{Spec}\mathbb Z[\mu_p,1/p],\mu_p^{\otimes2})
 \simeq (\operatorname{Cl}(F_1)/p)\otimes\mu_p
\]
is generated by cups of the $p$-units $1-\zeta_p^a$.
The class-group character in this component is
$\chi=\omega^{-1}\theta$; thus this is exactly the coefficient twist
required for the unit pairing. The excluded characters
$\theta=1,\omega^2$ correspond to $\chi=\omega^{-1},\omega$.
Their class-group components are zero by Stickelberger's theorem;
this exclusion is also explicitly checked in
\cite[proof of Lemma~5.2.7, p.~13]{WakeWangErickson2017}.
The remaining odd $\theta$-components of the displayed target vanish
because Vandiver gives the zero plus part of the class group.
We have therefore proved surjectivity of the full $p$-unit cup pairing.

The parity argument already proved above converts this into precisely
the hypothesis of Theorem~\ref{thm:pole-general}: the minus part of
the $p$-units modulo $p$ is the one-dimensional root-of-unity line,
its alternating self-cup is zero, and mixed-parity cups have zero
target. Thus plus Kummer classes suffice, and untwisting
$\mu_p=\mathbb F_p(1)$ makes them the odd-character part of
$H^1(P,\mathbb F_p)$. Theorem~\ref{thm:pole-general} gives the
kernel-field equality.

For $p=3$, the class group of $\mathbb Q(\mu_3)$ is zero, so the
same Kummer and Brauer calculation gives $H^2(P,\mathbb F_3)=0$.
The cup condition is vacuous and Theorem~\ref{thm:pole-general}
again applies. Finally, the equivalence of the Vandiver hypothesis
with the stated positive $K_{4k}$ vanishings is the known comparison
(K2)--(K3), proved in Section~\ref{sec:k-groups} using
\cite[VI.8.2, VI.8.3 and VI.10.9]{Weibel2013} and
\cite[Corollary~1.5, Proposition~3.7 and Remark~3.11]{Kurihara1992}.
\end{proof}

\begin{corollary}[K17-KERNEL-LT2POWER31]\label{cor:kernel-lt2power31}
For every prime $p<2^{31}$, one has $\Omega_p^*=\Omega_p$.
\end{corollary}

\begin{proof}
Hart--Harvey--Ong verify Vandiver for every odd prime below $2^{31}$
\cite[\S1, item~(2), p.~1]{HartHarveyOng2017}; apply
Theorem~\ref{thm:vandiver-kernel}. The prime $2$ is supplied by
Theorem~\ref{thm:pro2-kernel}.
\end{proof}

This numerical corollary uses the cited published computation. It
requires neither a new modular-symbol calculation nor any condition
on the multiplicities of the Iwasawa roots. For primes beyond this
range the same deduction applies whenever Vandiver is established.
The unconditional all-prime statement remains open in this argument.
